\documentclass[12pt,a4paper]{exam}
\usepackage[utf8]{inputenc}
\usepackage{amsmath,amssymb,amsthm}
\usepackage[margin=1in]{geometry}
\usepackage{graphicx}
\usepackage[colorlinks=true]{hyperref}
\pagestyle{headandfoot}
\firstpageheader{MIT OpenCourseWare}{}{\textbf{EXTREMELY HARD -- MIT LEVEL}}
\runningheader{MIT OpenCourseWare}{}{\textbf{EXTREMELY HARD -- MIT LEVEL}}
\firstpagefooter{}{Page \thepage}{}
\runningfooter{}{Page \thepage}{}

\begin{document}

\begin{center}
{\LARGE \textbf{MIT OpenCourseWare}}\\18.100A\\Real Analysis
\vspace{0.5cm}
{18.100A} --- {Real Analysis}
\vspace{0.3cm}
May 2026
\vspace{0.3cm}
\textbf{EXTREMELY HARD -- MIT LEVEL}
\end{center}

\begin{center}
\textbf{Time Limit: 3 hours}
\end{center}

\begin{center}
\textbf{Total Marks: 100}
\end{center}

\vspace{0.5cm}

\begin{questions}

\question[5]
Prove using the epsilon-delta definition that the function $f(x) = x^2 \sin(1/x)$ for $x \neq 0$ and $f(0) = 0$ is differentiable at $x = 0$ but $f'(x)$ is not continuous at $x = 0$. Carefully justify all steps in your epsilon-delta arguments.

\vspace{0.3cm}

\question[3]
Prove the Rellich-Kondrachov compactness theorem: $W^{1,p}(\Omega) \hookrightarrow\!\!\!\to L^q(\Omega)$ compactly for $1 \leq q < p^*$ when $\Omega$ is bounded and $\partial\Omega$ is $C^1$.

\vspace{0.3cm}

\question[2]
Let $M$ be a compact oriented $n$-manifold with boundary $\partial M$. State and prove Stokes' theorem for differential forms. Then apply it to prove the divergence theorem: $\int_M \operatorname{div} F \, dV = \int_{\partial M} F \cdot n \, dS$.

\vspace{0.3cm}

\question[4]
Prove the Poincar\'e inequality: for $u \in W_0^{1,p}(\Omega)$ with $\Omega \subset \mathbb{R}^n$ bounded, $\|u\|_{L^p} \leq C \|\nabla u\|_{L^p}$. Show that the constant $C$ depends on the diameter of $\Omega$ and give the optimal constant for $p=2$ on an interval.

\vspace{0.3cm}

\question[4]
Prove that $\mathbb{R}P^3$ is diffeomorphic to $SO(3)$ by constructing explicit maps and verifying they are smooth inverses. Use this to compute $\pi_1(SO(3))$.

\vspace{0.3cm}

\question[3]
Prove that symplectic integrators preserve the symplectic structure $\omega = \sum dq^i \wedge dp_i$ of Hamiltonian systems. Show that the Verlet (St\"ormer) method is symplectic and has energy error bounded over exponentially long times for integrable systems.

\vspace{0.3cm}

\question[6]
Construct a sequence of continuous functions $f_n: [0,1] \to \mathbb{R}$ that converges pointwise to $f(x) = 0$ but does NOT converge uniformly. Prove your construction works and explain why uniform convergence fails.

\vspace{0.3cm}

\question[3]
Prove the Hille-Yosida theorem: a linear operator $A$ generates a strongly continuous contraction semigroup on a Banach space $X$ iff $A$ is closed, densely defined, and all $\lambda > 0$ are in the resolvent set with $\|R(\lambda,A)\| \leq 1/\lambda$.

\vspace{0.3cm}

\question[4]
Compute the Fourier transform of the distribution $p.v.(1/x)$ and prove that $\widehat{p.v.(1/x)}(\xi) = -i\pi \, \operatorname{sgn}(\xi)$. Use this to solve the Hilbert transform equation $H(u) = f$.

\vspace{0.3cm}

\question[3]
Let $(X, d)$ be a complete metric space and $T: X \to X$ a contraction mapping with contraction factor $k < 1$. Prove the Banach fixed-point theorem: $T$ has a unique fixed point. Additionally, derive the error estimate $\|x_n - x^*\| \leq \frac{k^n}{1-k} \|x_1 - x_0\|$ where $x_{n+1} = T(x_n)$.

\vspace{0.3cm}

\question[5]
Prove the Schauder fixed-point theorem using Schauder estimates and the Arzela-Ascoli theorem: a compact operator $T: X \to X$ on a Banach space has a fixed point if it maps a closed convex set into itself.

\vspace{0.3cm}

\question[4]
Prove the Divergence theorem for smooth vector fields on compact manifolds with boundary and derive the Green's identities.

\vspace{0.3cm}

\question[3]
Prove that the Fourier transform is an isometry on $L^2(\mathbb{R}^n)$ (Plancherel's theorem) and extends uniquely from $L^1 \cap L^2$ to a unitary operator on $L^2$.

\vspace{0.3cm}

\question[4]
Prove that there is no continuous linear functional $L: \ell^\infty \to \mathbb{R}$ that extends the limit functional on $c$ (convergent sequences) and satisfies $\|L\| = 1$. (Hint: use the Hahn-Banach theorem and consider the Banach limit.)

\vspace{0.3cm}

\question[3]
Prove the Sobolev embedding theorem: for $p > n$, $W^{1,p}(\mathbb{R}^n) \hookrightarrow C^{0,\alpha}(\mathbb{R}^n)$ with $\alpha = 1 - n/p$. In particular, show that if $u \in W^{1,p}(\mathbb{R}^n)$ with $p > n$, then $u$ is (almost everywhere equal to) a H\"older continuous function.

\vspace{0.3cm}

\question[6]
Prove the Lax-Milgram theorem and use it to show existence and uniqueness of weak solutions to the Dirichlet problem $-\Delta u = f$ in $\Omega$, $u|_{\partial\Omega} = 0$ for $f \in L^2(\Omega)$.

\vspace{0.3cm}

\question[5]
Prove the Brouwer fixed-point theorem using differential topology: any continuous map $f: \overline{B}^n \to \overline{B}^n$ has a fixed point. Use the fact that there is no smooth retraction of the ball onto its boundary.

\vspace{0.3cm}

\question[5]
Prove the open mapping theorem: if $T: X \to Y$ is a bounded linear operator between Banach spaces and $T$ is surjective, then $T$ maps open sets to open sets. Use it to prove the inverse mapping theorem.

\vspace{0.3cm}

\question[4]
State and prove the spectral theorem for compact self-adjoint operators on a Hilbert space. Show that such an operator has a countable set of eigenvalues that accumulate only at zero, and the eigenvectors form an orthonormal basis.

\vspace{0.3cm}

\question[4]
Determine all values of $p \in \mathbb{R}$ for which the series $\sum_{n=2}^{\infty} \frac{1}{n^p \ln n}$ converges. Prove your answer using the integral test and comparison tests.

\vspace{0.3cm}

\question[3]
State and prove the Implicit Function Theorem. Apply it to show that the equation $F(x,y) = y^5 + y + x = 0$ can be solved for $y = g(x)$ near $(0,0)$, and compute $g'(0)$.

\vspace{0.3cm}

\question[6]
Construct an exhaustion function (proper smooth function bounded below) on any smooth manifold and use it to prove that every smooth manifold admits a complete Riemannian metric.

\vspace{0.3cm}

\question[6]
Prove the Gagliardo-Nirenberg-Sobolev inequality: for $1 \leq p < n$, there exists $C = C(n,p) > 0$ such that $\|u\|_{L^{p^*}(\mathbb{R}^n)} \leq C \|\nabla u\|_{L^p(\mathbb{R}^n)}$ for all $u \in C_c^1(\mathbb{R}^n)$, where $p^* = np/(n-p)$.

\vspace{0.3cm}

\question[3]
Prove the maximum principle for elliptic PDEs: if $Lu = -\Delta u + c(x) u \geq 0$ in $\Omega$ with $c \geq 0$, then $\max_{\overline{\Omega}} u \leq \max_{\partial\Omega} u^+$. Show how this implies uniqueness for the Dirichlet problem.

\vspace{0.3cm}

\question[2]
Prove Sard's theorem: if $f: U \subset \mathbb{R}^n \to \mathbb{R}^m$ is $C^k$ with $k > \max(n-m, 0)$, then the set of critical values of $f$ has Lebesgue measure zero in $\mathbb{R}^m$.

\vspace{0.3cm}

\end{questions}

\newpage
\section*{Solutions}

\textbf{Solution 1:}

We prove $f'(0) = 0$ by showing $|(f(h) - f(0))/h - 0| = |h \sin(1/h)| \leq |h|$. For any $\varepsilon > 0$, choose $\delta = \varepsilon$. Then $|h| < \delta$ implies $|h \sin(1/h)| < \varepsilon$, so $f'(0)=0$. For $x \neq 0$, $f'(x) = 2x \sin(1/x) - \cos(1/x)$. This oscillates wildly near 0 since $\cos(1/x)$ has no limit. Hence $f'$ is not continuous at 0.

\vspace{0.5cm}

\textbf{Solution 2:}

Take a bounded sequence $\{u_n\}$ in $W^{1,p}$. By the Sobolev-Gagliardo-Nirenberg inequality, it's bounded in $L^{p^*}$. Use the Frechet-Kolmogorov criterion for compactness in $L^q$: show $\|u_n(\cdot + h) - u_n\|_{L^q} \to 0$ uniformly as $h \to 0$ using difference quotients and the $W^{1,p}$ bound. Extract a subsequence convergent in $L^q$ by the Arzela-Ascoli type argument.

\vspace{0.5cm}

\textbf{Solution 3:}

Stokes: $\int_M d\omega = \int_{\partial M} \omega$ for an $(n-1)$-form $\omega$. Proof: Use partition of unity to reduce to $\mathbb{R}^n$ case, where it follows from the fundamental theorem of calculus and Fubini. For divergence theorem, let $\omega = \sum_i (-1)^{i-1} F_i \, dx^1 \wedge \cdots \wedge \widehat{dx^i} \wedge \cdots \wedge dx^n$. Then $d\omega = (\sum_i \partial F_i/\partial x^i) dx^1 \wedge \cdots \wedge dx^n = \operatorname{div} F \, dV$. The boundary term gives $\int_{\partial M} F \cdot n \, dS$ by the Hodge star duality.

\vspace{0.5cm}

\textbf{Solution 4:}

For $u \in C_c^\infty(\Omega)$, write $u(x) = \int_{-\infty}^{x_1} \partial_1 u(t, x'') dt$. By H\"older, $|u(x)|^p \leq (\operatorname{diam}\Omega)^{p-1} \int |\partial_1 u|^p dt$. Integrate over $\Omega$ to get $\|u\|_p^p \leq (\operatorname{diam}\Omega)^p \|\nabla u\|_p^p$. On $(0,\pi)$, the optimal constant for $p=2$ is $C = 1/\pi$, attained by $u(x) = \sin(x)$.

\vspace{0.5cm}

\textbf{Solution 5:}

Map $\mathbb{R}P^3 \to SO(3)$: represent a point in $\mathbb{R}P^3$ as a unit quaternion $q$ modulo sign. The conjugation map $v \mapsto q v q^{-1}$ on pure quaternions gives a rotation in $SO(3)$. This is a smooth double cover $S^3 \to SO(3)$ with kernel $\{\pm1\}$, inducing $\mathbb{R}P^3 \cong SO(3)$. Thus $\pi_1(SO(3)) \cong \mathbb{Z}/2\mathbb{Z}$.

\vspace{0.5cm}

\textbf{Solution 6:}

A one-step method $\Phi_h$ is symplectic if $(\Phi_h)^*\omega = \omega$. For Verlet: $q_{n+1/2} = q_n + \frac{h}{2} M^{-1} p_n$, $p_{n+1} = p_n - h \nabla V(q_{n+1/2})$, $q_{n+1} = q_{n+1/2} + \frac{h}{2} M^{-1} p_{n+1}$. Direct computation shows the Jacobian satisfies $J^T \Omega J = \Omega$. By the backward error analysis, the numerical solution is the exact flow of a nearby Hamiltonian, giving $O(e^{-c/h})$ energy conservation over exponentially long times.

\vspace{0.5cm}

\textbf{Solution 7:}

Let $f_n(x) = \max(0, 1 - n|x - 1/n|)$ (triangular spikes at $1/n$). For any $x$, for $n > 1/x$, $f_n(x) = 0$, so $f_n \to 0$ pointwise. But $\|f_n\|_\infty = 1$ for all $n$, so $\sup_{x \in [0,1]} |f_n(x) - 0| = 1 \not\to 0$, hence convergence is not uniform. This shows pointwise convergence does not imply uniform convergence.

\vspace{0.5cm}

\textbf{Solution 8:}

($\Rightarrow$): If $A$ generates $T(t)$, show $R(\lambda,A) = \int_0^\infty e^{-\lambda t} T(t) dt$, so $\|R(\lambda,A)\| \leq \int_0^\infty e^{-\lambda t} \|T(t)\| dt \leq 1/\lambda$. ($\Leftarrow$): Define $A_n = nAR(n,A)$ (Yosida approximants) which are bounded generators of $T_n(t) = e^{t A_n}$. Show $T_n(t)$ converges strongly to a semigroup $T(t)$ and its generator is $A$.

\vspace{0.5cm}

\textbf{Solution 9:}

$\langle \widehat{p.v.(1/x)}, \varphi \rangle = \langle p.v.(1/x), \widehat{\varphi} \rangle$. Write $p.v.(1/x) = \lim_{\varepsilon \to 0^+} \frac{1}{x} \chi_{|x|>\varepsilon}$. Its Fourier transform: $\int_{|x|>\varepsilon} \frac{e^{-2\pi i x \xi}}{x} dx = -i\pi \operatorname{sgn}(\xi)$ as $\varepsilon \to 0$. The Hilbert transform $H(u) = \frac{1}{\pi} p.v. \int \frac{u(t)}{x-t} dt$ has symbol $\widehat{H(u)}(\xi) = -i \operatorname{sgn}(\xi) \widehat{u}(\xi)$, so $H^{-1} = -H$.

\vspace{0.5cm}

\textbf{Solution 10:}

Pick any $x_0 \in X$, define $x_{n+1} = T(x_n)$. Then $d(x_{n+1}, x_n) \leq k^n d(x_1, x_0)$. By triangle inequality, $d(x_{n+m}, x_n) \leq \sum_{i=0}^{m-1} k^{n+i} d(x_1, x_0) \leq \frac{k^n}{1-k} d(x_1, x_0)$. So $\{x_n\}$ is Cauchy, converges to $x^*$. By continuity of $T$, $T(x^*) = x^*$. Uniqueness: if $Tx=x$, $Ty=y$, then $d(x,y) = d(Tx,Ty) \leq k d(x,y)$, so $d(x,y)=0$. The error estimate follows from taking $m \to \infty$.

\vspace{0.5cm}

\textbf{Solution 11:}

Let $C \subset X$ be closed convex, $T(C) \subset C$, $T$ compact. Approximate $T$ by finite-rank operators $T_n$ using Schauder basis. By Brouwer, each $T_n$ has a fixed point $x_n$. By compactness of $T$, $\{T(x_n)\}$ has convergent subsequence $T(x_{n_k}) \to y$. Since $x_{n_k} = T_{n_k}(x_{n_k})$ and $T_{n_k} \to T$ uniformly, $y$ is a fixed point of $T$.

\vspace{0.5cm}

\textbf{Solution 12:}

For $F \in \mathfrak{X}(\Omega)$, choose a partition of unity subordinated to a cover by coordinate charts that are half-balls near boundary. In each chart, the divergence theorem reduces to integrating a partial derivative over a cube (or half-cube), which by Fubini and FTC gives boundary integrals. Summing yields $\int_\Omega \operatorname{div} F \, dV = \int_{\partial\Omega} F \cdot n \, dS$. Green's identities follow by setting $F = u \nabla v$.

\vspace{0.5cm}

\textbf{Solution 13:}

For $f \in L^1 \cap L^2$, define $\widehat{f}(\xi) = \int f(x) e^{-2\pi i x \cdot \xi} dx$. Show $\int |\widehat{f}|^2 = \int |f|^2$ by proving $\int \widehat{f} \bar{g} = \int f \bar{g}$ for $g$ Schwartz (using Gaussian approximation and Fubini). Extend by density of $L^1 \cap L^2$ in $L^2$. The inverse Fourier transform gives the adjoint, and $\mathcal{F}^* = \mathcal{F}^{-1}$.

\vspace{0.5cm}

\textbf{Solution 14:}

The Hahn-Banach theorem guarantees the existence of such extensions (Banach limits), but they are not unique and not continuous in the sense of being representable by a sequence. The key issue: $\ell^\infty / c_0$ is non-separable and the quotient map is not injective. Banach limits exist but require the axiom of choice. They satisfy $\liminf x_n \leq L(x) \leq \limsup x_n$ and are shift-invariant.

\vspace{0.5cm}

\textbf{Solution 15:}

For $u \in C_c^\infty(\mathbb{R}^n)$, write $u(x) - u(y)$ as an integral of $\nabla u$ over a chain connecting $x$ and $y$. Using the mean value integral representation and H\"older inequality, $|u(x) - u(y)| \leq C \|\nabla u\|_{L^p} |x-y|^{1-n/p}$. By density of smooth functions in $W^{1,p}$, the inequality extends, proving the embedding is continuous.

\vspace{0.5cm}

\textbf{Solution 16:}

Lax-Milgram: Let $a: H \times H \to \mathbb{R}$ be coercive bilinear continuous on Hilbert $H$, $f \in H^*$. Then $\exists! u \in H$: $a(u,v) = f(v)$ for all $v \in H$. For the Dirichlet problem, $H = H_0^1(\Omega)$, $a(u,v) = \int_\Omega \nabla u \cdot \nabla v$, $f(v) = \int_\Omega f v$. Poincar\'e inequality gives coercivity, trace theorem gives boundary values.

\vspace{0.5cm}

\textbf{Solution 17:}

Suppose $f$ has no fixed point. Define $r(x) =$ intersection of ray from $f(x)$ through $x$ with $\partial B^n$. This gives a smooth retraction (after smooth approximation). By Sard's theorem, regular values exist. For a regular value $y$, $r^{-1}(y)$ is a 1-manifold with boundary, one endpoint at $y$. The other endpoint would be in $B^n$, impossible since $r$ maps interior to boundary. Contradiction.

\vspace{0.5cm}

\textbf{Solution 18:}

Use Baire category: $Y = \cup_n n T(B_X(0,1))$ so some $nT(B_X(0,1))$ has nonempty interior. Show $T(B_X(0,1))$ contains a neighborhood of 0. Then by scaling, $T$ maps open balls to open sets. The inverse mapping theorem follows: if $T$ is bijective bounded linear, then $T^{-1}$ is bounded.

\vspace{0.5cm}

\textbf{Solution 19:}

Let $T: H \to H$ be compact self-adjoint. Show $\|T\| = \sup_{|x|=1} |\langle Tx, x \rangle|$ (achieved). Let $\lambda_1 = \pm\|T\|$ be the extremal eigenvalue with eigenvector $e_1$. Inductively apply the same argument to $T$ restricted to $\{e_1,\ldots,e_{n-1}\}^\perp$. Compactness implies $\lambda_n \to 0$. The eigenvectors form an ONB by the spectral decomposition $Tx = \sum \lambda_n \langle x, e_n \rangle e_n$.

\vspace{0.5cm}

\textbf{Solution 20:}

By the integral test, $\int_2^{\infty} \frac{dx}{x^p \ln x}$ converges iff $p > 1$. For $p=1$, $\int_2^{\infty} \frac{dx}{x \ln x} = \lim_{b \to \infty} \ln(\ln b) - \ln(\ln 2) = \infty$, so diverges. For $p<1$, compare with $1/x^p$: $\frac{1}{n^p \ln n} > \frac{1}{n^p \cdot n^{\varepsilon}}$ for large $n$, which diverges by p-test. For $p>1$, $\frac{1}{n^p \ln n} < \frac{1}{n^{(p+1)/2}}$ for large $n$, which converges.

\vspace{0.5cm}

\textbf{Solution 21:}

IFT: Let $F: \mathbb{R}^{m+n} \to \mathbb{R}^n$ be $C^1$, $F(a,b)=0$, and $\partial F/\partial y (a,b)$ invertible. Then there exists $g: U \to \mathbb{R}^n$ $C^1$ near $a$ with $F(x,g(x))=0$ and $Dg(x) = -(\partial F/\partial y)^{-1} \partial F/\partial x$. For $F(x,y) = y^5 + y + x$, $\partial F/\partial y = 5y^4 + 1$, at $(0,0)$ this is $1 \neq 0$, so IFT applies. $g'(0) = - (\partial F/\partial y)^{-1} \partial F/\partial x = -1/1 = -1$.

\vspace{0.5cm}

\textbf{Solution 22:}

Let $\{U_i\}$ be a countable cover by coordinate balls, $\{\rho_i\}$ a partition of unity. Define $f = \sum_i i \rho_i$. This is proper: $f^{-1}([-\infty, c])$ is compact (only finitely many $i \leq c$ contribute). For a Riemannian metric, start with any metric $g_0$, let $g = g_0 + df \otimes df$ on $f^{-1}((-\infty, \infty))$. The gradient of $f$ has unit length, so any curve must escape to infinity in infinite length.

\vspace{0.5cm}

\textbf{Solution 23:}

For $p=1$, use $|u(x)| \leq \int_{-\infty}^{x_i} |\partial_i u| dx_i$, multiply over $i=1,\ldots,n$, integrate, apply H\"older repeatedly to get $\|u\|_{n/(n-1)} \leq \prod (\int |\partial_i u|)^{1/n} \leq \frac{1}{n} \sum \int |\partial_i u|$. For general $p$, apply the $p=1$ case to $|u|^q$ with suitable $q$ and use H\"older.

\vspace{0.5cm}

\textbf{Solution 24:}

Suppose $u$ attains an interior maximum at $x_0$ with $u(x_0) = M > 0$ and $M > \max_{\partial\Omega} u^+$. At $x_0$, $\Delta u \leq 0$ (Hessian negative semidefinite), so $Lu = -\Delta u + c u \geq 0 + cM \geq 0$. Strictly, if $c>0$ or $M>0$, this gives $Lu > 0$ at the maximum, contradiction. For uniqueness of $\Delta u = f$ with $u|_{\partial\Omega} = g$, difference of two solutions satisfies $\Delta w = 0$, $w|_{\partial\Omega} = 0$, so $w \equiv 0$.

\vspace{0.5cm}

\textbf{Solution 25:}

Cover $U$ by cubes. On each cube where $\|Df\| \leq C$ and $|\det Df| \leq \varepsilon$ on a critical point, use the Whitney extension to show the image can be covered by $O((\varepsilon/C)^{m})$ balls of radius $C/\sqrt{n}$. Summing over cubes gives arbitrarily small total measure. The differentiability condition ensures the critical set has enough structure.

\vspace{0.5cm}

\end{document}